
\documentclass[11pt]{article}
\usepackage[spanish]{babel}
\usepackage[utf8]{inputenc}
\usepackage[T1]{fontenc}
\usepackage[margin=2cm]{geometry}
\usepackage{amsmath,amssymb,amsthm}
\usepackage{enumitem}
\usepackage[hidelinks]{hyperref}
\usepackage{longtable}
\usepackage{array}
\usepackage{booktabs}
\usepackage{tikz}
\usepackage{xcolor}

\setlength{\parindent}{0pt}
\setlength{\parskip}{0.75\baselineskip}

\newtheoremstyle{prismstyle}
  {0.8\baselineskip}
  {0.8\baselineskip}
  {}
  {}
  {\bfseries}
  {.}
  {0.5em}
  {}

\theoremstyle{prismstyle}
\newtheorem{problem}{Problema}

\newenvironment{solution}[1][Solución]
  {\par\textbf{#1.}\ }
  {\par}

\newcommand{\problemsection}[2]{
  \section{#1}
  \begin{problem}
  #2
  \end{problem}
}

\newcolumntype{L}[1]{>{\raggedright\arraybackslash}p{#1}}

\title{Problemas de Cónicas}
\author{}
\date{\today}

\begin{document}
\maketitle

\tableofcontents


\problemsection{Problema 1: Otra definición de una elipse}{
Una elipse se puede definir como el lugar geométrico de los puntos cuya suma de las distancias a dos puntos fijos llamados focos es constante; esto es, si $P$ es un punto de la elipse y $F_1$, $F_2$ son sus focos, entonces
\[
F_1P + F_2P = \text{constante}.
\]

A partir de esta definición y de los parámetros descritos previamente, demuestre lo siguiente:

\begin{enumerate}[label=\alph*)]
  \item La constante a la que es igual dicha suma es $2a$.
  \item Los parámetros $a$, $b$ y la distancia focal $c$ satisfacen la relación pitagórica
  \[
  a^2 = b^2 + c^2.
  \]
  \item A partir de esta definición, la ecuación de la elipse horizontal centrada en $C=(0,0)$ es
  \[
  \frac{x^2}{a^2} + \frac{y^2}{b^2} = 1.
  \]
  \item La excentricidad también puede calcularse como
  \[
  e = \frac{c}{a}.
  \]
  \item El semilatus rectum también puede calcularse como
  \[
  p = \frac{b^2}{a}.
  \]
\end{enumerate}
}

\begin{solution}
\textbf{Inciso a) Demostración de que la constante es $2a$.}

Colocamos la elipse en un sistema cartesiano con centro en $C(0,0)$ y focos en $F_1(-c,0)$ y $F_2(c,0)$. Los vértices sobre el eje mayor están en $V_1(-a,0)$ y $V_2(a,0)$.

Por definición, para cualquier punto $P$ de la elipse se cumple que $F_1P + F_2P = K$, donde $K$ es constante. Evaluamos la suma en el vértice derecho $V_2(a,0)$:
\[
F_1V_2 = \text{distancia}((-c,0),(a,0)) = a-(-c) = a+c,
\]
\[
F_2V_2 = \text{distancia}((c,0),(a,0)) = a-c.
\]
Por tanto,
\[
K = (a+c) + (a-c) = 2a.
\]
Como la suma es constante para toda la curva, concluimos que
\[
\boxed{F_1P + F_2P = 2a}.
\]

\textbf{Inciso b) Demostración de $a^2 = b^2 + c^2$.}

Consideremos el punto $B(0,b)$, extremo del semieje menor. Por simetría, las distancias a los focos son iguales, y además
\[
F_1B = \sqrt{(-c-0)^2 + (0-b)^2} = \sqrt{c^2+b^2},
\]
\[
F_2B = \sqrt{(c-0)^2 + (0-b)^2} = \sqrt{c^2+b^2}.
\]
Usando el resultado del inciso anterior,
\[
F_1B + F_2B = 2\sqrt{b^2+c^2} = 2a.
\]
Dividiendo entre $2$ y elevando al cuadrado se obtiene
\[
\sqrt{b^2+c^2} = a \quad \Rightarrow \quad \boxed{a^2 = b^2+c^2}.
\]

\textbf{Inciso c) Derivación de la ecuación canónica.}

Sea $P(x,y)$ un punto genérico de la elipse. Entonces,
\[
\sqrt{(x+c)^2 + y^2} + \sqrt{(x-c)^2 + y^2} = 2a.
\]
Despejamos uno de los radicales:
\[
\sqrt{(x+c)^2 + y^2} = 2a - \sqrt{(x-c)^2 + y^2}.
\]
Elevando al cuadrado ambos lados,
\[
(x+c)^2 + y^2 = 4a^2 - 4a\sqrt{(x-c)^2 + y^2} + (x-c)^2 + y^2.
\]
Al expandir y simplificar,
\[
4cx = 4a^2 - 4a\sqrt{(x-c)^2 + y^2}.
\]
Dividimos entre $4$ y reordenamos:
\[
a^2 - cx = a\sqrt{(x-c)^2 + y^2}.
\]
Volvemos a elevar al cuadrado:
\[
(a^2-cx)^2 = a^2\bigl((x-c)^2+y^2\bigr).
\]
Desarrollando ambos lados,
\[
a^4 - 2a^2cx + c^2x^2 = a^2x^2 - 2a^2cx + a^2c^2 + a^2y^2.
\]
Cancelando términos comunes y reagrupando,
\[
a^2(a^2-c^2) = (a^2-c^2)x^2 + a^2y^2.
\]
Como por el inciso b) se tiene $a^2-c^2=b^2$, resulta
\[
a^2b^2 = b^2x^2 + a^2y^2.
\]
Dividiendo entre $a^2b^2$, obtenemos la ecuación canónica:
\[
\boxed{\frac{x^2}{a^2} + \frac{y^2}{b^2} = 1}.
\]

\textbf{Inciso d) Demostración de $e = c/a$.}

Del inciso b) se sigue que
\[
c^2 = a^2-b^2.
\]
Por otra parte, en la parametrización usual de la elipse,
\[
b^2 = a^2(1-e^2).
\]
Sustituyendo en la expresión de $c^2$,
\[
c^2 = a^2 - a^2(1-e^2) = a^2e^2.
\]
Tomando la raíz positiva,
\[
c = ae \quad \Rightarrow \quad \boxed{e = \frac{c}{a}}.
\]

\textbf{Inciso e) Demostración de $p = b^2/a$.}

El semilatus rectum $p$ es la distancia perpendicular desde un foco hasta la curva; analíticamente corresponde al valor de $|y|$ cuando $x=c$.
Sustituyendo $x=c$ en la ecuación canónica,
\[
\frac{c^2}{a^2} + \frac{y^2}{b^2} = 1.
\]
Entonces,
\[
\frac{y^2}{b^2} = 1 - \frac{c^2}{a^2} = \frac{a^2-c^2}{a^2}.
\]
Usando nuevamente que $a^2-c^2=b^2$,
\[
\frac{y^2}{b^2} = \frac{b^2}{a^2} \quad \Rightarrow \quad y^2 = \frac{b^4}{a^2}.
\]
Tomando la raíz positiva,
\[
\boxed{p = |y| = \frac{b^2}{a}}.
\]
\end{solution}

\problemsection{Problema 2: Ecuación de una elipse}{
Encuentre la ecuación general de la elipse cuyo semilatus rectum valga $9$ y cuya excentricidad sea $0.5$.
}

\begin{solution}
Para resolver este problema, partimos de las relaciones geométricas fundamentales de una elipse en su posición canónica, es decir, centrada en el origen y con eje mayor sobre el eje $x$. Adoptamos esta configuración estándar porque el enunciado no indica traslaciones ni rotaciones.

\textbf{Paso 1. Datos conocidos y relaciones básicas.}

Se nos proporcionan los parámetros
\[
p = 9, \qquad e = 0.5.
\]
En una elipse horizontal se cumplen las relaciones
\[
p = a(1-e^2)
\qquad \text{y} \qquad
b^2 = a^2(1-e^2).
\]
De ellas también se deduce la identidad útil
\[
b^2 = ap.
\]

\textbf{Paso 2. Cálculo del semieje mayor $a$.}

Despejamos $a$ de la fórmula del semilatus rectum:
\[
a = \frac{p}{1-e^2} = \frac{9}{1-(0.5)^2} = \frac{9}{0.75} = 12.
\]

\textbf{Paso 3. Cálculo del semieje menor.}

Usando la relación $b^2 = ap$,
\[
b^2 = 12 \cdot 9 = 108.
\]
Así, los parámetros principales de la elipse son
\[
a^2 = 144, \qquad b^2 = 108.
\]

\textbf{Paso 4. Ecuación canónica.}

Sustituyendo en la forma estándar de la elipse horizontal,
\[
\frac{x^2}{a^2} + \frac{y^2}{b^2} = 1,
\]
obtenemos
\[
\frac{x^2}{144} + \frac{y^2}{108} = 1.
\]

\textbf{Paso 5. Conversión a la ecuación general.}

Multiplicamos toda la ecuación por el mínimo común múltiplo de $144$ y $108$, que es $432$:
\[
432\left(\frac{x^2}{144}\right) + 432\left(\frac{y^2}{108}\right) = 432.
\]
Entonces,
\[
3x^2 + 4y^2 = 432.
\]
Llevando todos los términos al mismo lado,
\[
\boxed{3x^2 + 4y^2 - 432 = 0}.
\]

\textbf{Verificación.}

Comprobamos que la ecuación obtenida reproduce los datos del problema:
\[
e = \sqrt{1-\frac{b^2}{a^2}} = \sqrt{1-\frac{108}{144}} = \sqrt{0.25} = 0.5,
\]
y además
\[
p = \frac{b^2}{a} = \frac{108}{12} = 9.
\]
Ambos valores coinciden con el enunciado, así que la ecuación es correcta.
\end{solution}

\problemsection{Problema 3: Parámetros geométricos}{
Defina conceptualmente a qué corresponden cada uno de los siguientes parámetros en el caso de una elipse:
\begin{enumerate}[label=\alph*)]
  \item defina $F$
  \item defina $e$
  \item defina $p$
  \item defina $C$
  \item defina $a$
  \item defina $b$
  \item defina $r$
  \item defina $f$
  \item defina $E$
\end{enumerate}
}

\begin{solution}
A continuación se presenta la definición conceptual de cada parámetro, estructurada en una tabla que integra su significado geométrico, su relación matemática principal y su interpretación en mecánica celeste.

\setlength{\LTleft}{0pt}
\setlength{\LTright}{0pt}
\setlength{\tabcolsep}{3pt}
\renewcommand{\arraystretch}{1.18}
{\footnotesize
\begin{longtable}{@{}L{0.11\textwidth}L{0.08\textwidth}L{0.23\textwidth}L{0.24\textwidth}L{0.22\textwidth}@{}}
\toprule
\textbf{Parámetro} & \textbf{Símbolo} & \textbf{Definición conceptual} & \textbf{Relación matemática / propiedad geométrica} & \textbf{Interpretación en mecánica celeste} \\
\midrule
\endfirsthead
\toprule
\textbf{Parámetro} & \textbf{Símbolo} & \textbf{Definición conceptual} & \textbf{Relación matemática / propiedad geométrica} & \textbf{Interpretación en mecánica celeste} \\
\midrule
\endhead
\bottomrule
\endfoot
\textbf{Foco} & $F$ & Punto fijo del plano respecto al cual se construye la elipse mediante la propiedad de suma constante de distancias. Es el origen del sistema de coordenadas polares perifocales. & Distancia al centro: $c = ae$. Existen dos focos simétricos respecto a $C$. En el problema de dos cuerpos, el foco activo coincide con la posición del cuerpo masivo. & Centro de atracción gravitacional. El vector posición $\vec{r}$ se mide desde este punto. \\
\textbf{Excentrici-}\textbf{dad} & $e$ & Parámetro adimensional que cuantifica el grado de achatamiento o desviación de la elipse respecto a una circunferencia. & $e = \dfrac{c}{a} = \sqrt{1-\dfrac{b^2}{a^2}}$, con $0 \le e < 1$ para elipses. Si $e=0$, la elipse degenera en circunferencia. & Clasifica el tipo de cónica y determina la variación radial entre periapsis y apoapsis. \\
\textbf{Semilatus rectum} & $p$ & Distancia perpendicular desde un foco hasta la curva, medida a lo largo de una recta paralela al eje menor. Caracteriza la apertura de la cónica en la vecindad del foco. & $p = a(1-e^2) = \dfrac{b^2}{a}$. Aparece en la ecuación polar $r(\theta) = \dfrac{p}{1+e\cos f}$. & Se relaciona con el momento angular específico $h$ y el parámetro gravitacional $\mu$ mediante $p = h^2/\mu$. \\
\textbf{Centro geométrico} & $C$ & Punto de intersección de los ejes de simetría mayor y menor. Es el centro de simetría puntual de la figura. & En coordenadas cartesianas canónicas, $(x_C,y_C)=(0,0)$. La distancia a cada foco es $c = ae$. & Referencia para la descripción cartesiana de la elipse. No coincide con el centro de fuerzas gravitacionales. \\
\textbf{Semieje mayor} & $a$ & Mitad de la longitud máxima de la elipse a lo largo de su eje principal. Representa la escala característica de tamaño de la órbita. & $a = \dfrac{p}{1-e^2} = \dfrac{r_{\text{peri}} + r_{\text{apo}}}{2}$ y satisface $a^2 = b^2 + c^2$. & Parámetro fundamental en la tercera ley de Kepler: $n^2 a^3 = \mu$. También define la energía específica $\mathcal{E} = -\mu/(2a)$. \\
\textbf{Semieje menor} & $b$ & Mitad de la longitud máxima a lo largo del eje secundario, perpendicular al eje mayor. Cuantifica la extensión transversal de la órbita. & $b = a\sqrt{1-e^2}$. Corresponde a la coordenada $y$ cuando $x=0$ en el sistema centrado en $C$. & Determina el área total de la elipse, $A = \pi ab$, y se vincula con la velocidad areal constante. \\
\textbf{Radio vector} & $r$ & Distancia instantánea desde el foco $F$ hasta un punto arbitrario sobre la elipse. Es la magnitud del vector posición $\vec{r}$. & $r = \|\vec{r}\| = a(1-e\cos E) = \dfrac{p}{1+e\cos f}$, con $a(1-e) \le r \le a(1+e)$. & Evoluciona en el tiempo según la segunda ley de Kepler y su variación está acotada por los ápsides. \\
\textbf{Anomalía verdadera} & $f$ & Ángulo polar medido desde la dirección del periapsis hasta el vector posición $\vec{r}$. Indica la posición angular real del cuerpo. & $0 \le f < 2\pi$. Se relaciona con $E$ mediante $\tan\dfrac{f}{2} = \sqrt{\dfrac{1+e}{1-e}}\tan\dfrac{E}{2}$. & Es el parámetro natural para la ecuación polar y describe la geometría instantánea observada desde el foco. \\
\textbf{Anomalía excéntrica} & $E$ & Parámetro angular auxiliar definido geométricamente proyectando el punto de la elipse sobre una circunferencia auxiliar de radio $a$ centrada en $C$. & $0 \le E < 2\pi$. Vincula tiempo y posición mediante la ecuación de Kepler: $M = E - e\sin E$. & Variable intermedia esencial para resolver el problema de Kepler en función del tiempo. \\
\end{longtable}
}

\textbf{Análisis complementario y relaciones estructurales.}

Los nueve parámetros enumerados no son independientes; forman un sistema algebraico y geométrico cohesionado que permite transitar entre descripciones cartesianas, polares y temporales de la elipse.

\begin{enumerate}
  \item \textbf{Tríada geométrica básica.} Los parámetros $a$, $b$ y $c$ satisfacen la relación
  \[
  a^2 = b^2 + c^2,
  \]
  donde $c = ae$ es la distancia entre el foco y el centro.

  \item \textbf{Transición polar-cartesiana.} El radio vector $r$ y la anomalía verdadera $f$ determinan completamente la posición en coordenadas polares focales. Para integrar la dinámica temporal se introduce la anomalía excéntrica $E$, con la cual el área barrida se expresa como
  \[
  \Delta A = \frac{ab}{2}(E - e\sin E).
  \]
  Al compararla con la ley de áreas, se obtiene la ecuación de Kepler.

  \item \textbf{Invariancia del semilatus rectum.} El parámetro $p$ encapsula la conservación del momento angular específico, pues
  \[
  p = \frac{h^2}{\mu}.
  \]

  \item \textbf{Jerarquía de dependencia.} En la práctica teórica y computacional, los parámetros pueden organizarse en niveles: geométrico primario $(a,e)$; derivado directo $b = a\sqrt{1-e^2}$, $c = ae$, $p = a(1-e^2)$; cinemático $r(f)$, $r(E)$, $f(E)$, $E(t)$; y energético $\mathcal{E} = -\mu/(2a)$, $h = \sqrt{\mu p}$.
\end{enumerate}

Esta estructura paramétrica permite reconstruir la geometría orbital y predecir la posición futura del sistema a partir de un conjunto reducido de parámetros fundamentales.
\end{solution}

\problemsection{Problema 6: Órbitas de cometas}{
El semieje mayor de un cometa es de $20\ \text{UA}$ y la distancia entre el foco y el perihelio es de $1\ \text{UA}$.
\begin{enumerate}[label=\alph*)]
  \item ¿Cuál es la excentricidad de la órbita?
  \item ¿Cuánto vale el semieje menor?
  \item ¿Cuánto vale el semilatus rectum?
  \item Escriba la ecuación de este cometa en coordenadas polares.
  \item Haga un bosquejo de la situación medianamente a escala.
\end{enumerate}
}

\begin{solution}
\textbf{Datos iniciales.}
\[
a = 20\ \text{UA},
\qquad
q = 1\ \text{UA},
\]
donde $q$ representa la distancia foco-perihelio, es decir, el radio periapsidal.

\textbf{Inciso a) Cálculo de la excentricidad $e$.}

En una elipse, la distancia desde el foco hasta el perihelio satisface
\[
q = a(1-e).
\]
Despejando la excentricidad,
\[
e = 1 - \frac{q}{a} = 1 - \frac{1}{20} = 0.95.
\]
Por tanto,
\[
\boxed{e = 0.95}.
\]
Este valor, cercano a $1$, indica una órbita muy elongada, como es típico en muchos cometas.

\textbf{Inciso b) Cálculo del semieje menor $b$.}

Usamos la relación
\[
b = a\sqrt{1-e^2}.
\]
Sustituyendo los datos,
\[
b = 20\sqrt{1-(0.95)^2} = 20\sqrt{1-0.9025} = 20\sqrt{0.0975}.
\]
Numéricamente,
\[
b \approx 20(0.31225) \approx 6.245\ \text{UA}.
\]
En consecuencia,
\[
\boxed{b \approx 6.245\ \text{UA}}.
\]

\textbf{Inciso c) Cálculo del semilatus rectum $p$.}

El semilatus rectum se relaciona con $a$ y $e$ mediante
\[
p = a(1-e^2).
\]
Entonces,
\[
p = 20(1-0.9025) = 20(0.0975) = 1.95\ \text{UA}.
\]
Por tanto,
\[
\boxed{p = 1.95\ \text{UA}}.
\]
Como verificación adicional, también se cumple $p = q(1+e)$, y en efecto
\[
p = 1(1+0.95) = 1.95\ \text{UA}.
\]

\textbf{Inciso d) Ecuación en coordenadas polares.}

Con el foco en el origen y el eje polar apuntando hacia el perihelio, la ecuación polar de la órbita es
\[
r(\theta) = \frac{p}{1+e\cos\theta}.
\]
Sustituyendo los valores obtenidos,
\[
\boxed{r(\theta) = \frac{1.95}{1+0.95\cos\theta}}.
\]
Aquí $r$ está medido en unidades astronómicas y $\theta$ representa la anomalía verdadera medida desde la dirección del perihelio.

\textbf{Inciso e) Bosquejo medianamente a escala.}

Como $c = ae = 19\ \text{UA}$, el centro de la elipse queda en $(19,0)$ si ubicamos el foco en el origen y el perihelio en $(1,0)$. Además, el afelio queda en $(39,0)$ y el semieje menor vale aproximadamente $6.245\ \text{UA}$. El siguiente esquema respeta esas proporciones de manera aproximada:

\begin{center}
\shorthandoff{<>}
\begin{tikzpicture}[x=0.18cm,y=0.18cm]
  \draw[->] (-2,0) -- (42,0) node[right] {$x$};
  \draw[->] (0,-9) -- (0,9) node[above] {$y$};
  \draw[thick,blue] (19,0) ellipse [x radius=20, y radius=6.245];
  \fill (0,0) circle (0.28) node[below left] {$F$};
  \fill (19,0) circle (0.22) node[below] {$C$};
  \fill (1,0) circle (0.22) node[below] {Perihelio};
  \fill (39,0) circle (0.22) node[below] {Afelio};
  \draw[dashed] (19,-6.245) -- (19,6.245);
  \node[right] at (19,6.245) {$b \approx 6.245$};
  \draw[<->] (0,-7.5) -- (19,-7.5) node[midway,below] {$c=19$};
  \draw[<->] (19,-8.5) -- (39,-8.5) node[midway,below] {$a=20$};
\end{tikzpicture}
\shorthandon{<>}
\end{center}
\end{solution}

\problemsection{Problema 7: Órbita lunar}{
Sobre una órbita lunar, el punto más cercano al satélite se conoce como \textbf{periselenio} y el más lejano como \textbf{aposelenio}. La nave Apolo 11 fue puesta en una órbita elíptica alrededor de la Luna con una altitud respecto a la superficie lunar de $110\ \text{km}$ en el periselenio y de $314\ \text{km}$ en el aposelenio. Encuentre la ecuación de esta elipse en coordenadas polares si el radio de la Luna es de $1728\ \text{km}$ y el centro de la Luna se encuentra en uno de los focos de la elipse. Haga un bosquejo de la situación medianamente a escala.
}

\begin{solution}
\textbf{1. Datos del problema.}

Se tienen los siguientes datos:
\[
R_L = 1728\ \text{km},
\qquad
h_p = 110\ \text{km},
\qquad
h_a = 314\ \text{km}.
\]

\textbf{2. Distancias desde el centro de la Luna.}

Como el centro de la Luna ocupa uno de los focos de la elipse, las distancias radiales en periselenio y aposelenio son
\[
r_p = R_L + h_p = 1728 + 110 = 1838\ \text{km},
\]
\[
r_a = R_L + h_a = 1728 + 314 = 2042\ \text{km}.
\]

\textbf{3. Ecuación polar de la cónica.}

La ecuación de una cónica en coordenadas polares con origen en el foco es
\[
r = \frac{p}{1+e\cos f},
\]
donde $p$ es el semilatus rectum, $e$ la excentricidad y $f$ la anomalía verdadera medida desde el periselenio.

\textbf{4. Cálculo de la excentricidad $e$.}

Evaluamos la ecuación polar en los dos extremos de la órbita:
\[
r_p = \frac{p}{1+e}
\qquad (f=0),
\]
\[
r_a = \frac{p}{1-e}
\qquad (f=\pi).
\]
Despejando $p$ de ambas expresiones e igualando,
\[
r_p(1+e) = r_a(1-e).
\]
Desarrollando,
\[
r_p + r_pe = r_a - r_ae,
\]
\[
e(r_p+r_a) = r_a-r_p.
\]
Por tanto,
\[
\boxed{e = \frac{r_a-r_p}{r_a+r_p}}.
\]
Sustituyendo valores numéricos,
\[
e = \frac{2042-1838}{2042+1838} = \frac{204}{3880} = 0{,}052577 \approx \boxed{0{,}0526}.
\]

\textbf{5. Cálculo del semilatus rectum $p$.}

Usando la expresión del periselenio,
\[
p = r_p(1+e) = 1838(1+0{,}052577) = 1838(1{,}052577) \approx \boxed{1934{,}64\ \text{km}}.
\]
Como verificación,
\[
p = r_a(1-e) = 2042(1-0{,}052577) = 2042(0{,}947423) \approx 1934{,}64\ \text{km}.
\]

\textbf{6. Semieje mayor $a$.}

Para una elipse,
\[
a = \frac{r_p+r_a}{2} = \frac{1838+2042}{2} = \boxed{1940\ \text{km}}.
\]
Además, se verifica la consistencia con
\[
p = a(1-e^2).
\]
En efecto,
\[
p = 1940\left(1-(0{,}052577)^2\right) \approx 1940(0{,}997236) \approx 1934{,}64\ \text{km}.
\]

\textbf{7. Ecuación final en coordenadas polares.}

Sustituyendo los valores obtenidos en la ecuación polar,
\[
\boxed{r(f) = \frac{1934{,}64\ \text{km}}{1+0{,}0526\cos f}}.
\]
Aquí $f$ se mide desde la dirección del periselenio.

\textbf{8. Parámetros geométricos adicionales para el bosquejo.}

El semieje menor es
\[
b = a\sqrt{1-e^2} = 1940\sqrt{1-0{,}002764} \approx 1937{,}3\ \text{km}.
\]
La distancia foco-centro resulta
\[
c = ae = 1940(0{,}0526) \approx 102{,}0\ \text{km}.
\]

\textbf{9. Bosquejo medianamente a escala.}

Tomando el foco lunar en el origen, el periselenio queda en $(1838,0)$ y el aposelenio en $(-2042,0)$. El centro geométrico de la órbita se ubica aproximadamente en $(-102,0)$ y la Luna puede representarse como un círculo de radio $1728\ \text{km}$ centrado en el foco. El siguiente esquema refleja esta geometría de manera aproximada:

\begin{center}
\shorthandoff{<>}
\begin{tikzpicture}[x=0.0018cm,y=0.0018cm]
  \draw[->] (-2600,0) -- (2400,0) node[right] {$x$};
  \draw[->] (0,-2300) -- (0,2300) node[above] {$y$};
  \draw[fill=gray!15] (0,0) circle[radius=1728];
  \draw[thick,blue] (-102,0) ellipse [x radius=1940, y radius=1937.3];
  \fill (0,0) circle (60) node[above right] {Centro lunar / foco};
  \fill (1838,0) circle (40) node[below right] {Periselenio};
  \fill (-2042,0) circle (40) node[below left] {Aposelenio};
  \fill (-102,0) circle (35) node[above] {$C$};
\end{tikzpicture}
\shorthandon{<>}
\end{center}
\end{solution}

\problemsection{Problema 8: Anomalías}{
Demuestre que, para una elipse:
\begin{enumerate}[label=\alph*)]
  \item la distancia radial $r$ está relacionada con la anomalía excéntrica $E$ por medio de la expresión
  \[
  r = a(1 - e\cos E),
  \]
  \item y que la anomalía verdadera se relaciona con la anomalía excéntrica a través de
  \[
  \tan\frac{f}{2} = \sqrt{\frac{1+e}{1-e}}\tan\frac{E}{2}.
  \]
\end{enumerate}
}

\begin{solution}
\textbf{Inciso a) Demostración de $r = a(1 - e\cos E)$.}

\textbf{Paso 1. Configuración geométrica y parametrización.}

Colocamos el centro de la elipse en el origen de un sistema cartesiano $(x_c,y_c)$. Los focos se ubican en $F_1(-ae,0)$ y $F_2(ae,0)$, donde $a$ es el semieje mayor, $e$ la excentricidad y $c=ae$ la distancia del centro a cada foco. Además,
\[
b^2 = a^2(1-e^2).
\]

La elipse puede parametrizarse mediante la circunferencia auxiliar de radio $a$. Un punto $P$ sobre la elipse tiene coordenadas
\[
x_c = a\cos E,
\qquad
y_c = b\sin E,
\]
donde $E$ es la anomalía excéntrica.

\textbf{Paso 2. Expresión de la distancia radial $r$.}

Tomamos como origen polar el foco derecho $F_2$, ubicado en $x=ae$. Entonces la distancia radial desde el foco hasta el punto $P$ es
\[
r = \sqrt{(x_c-ae)^2 + y_c^2}.
\]
Sustituyendo la parametrización,
\[
r^2 = (a\cos E-ae)^2 + (b\sin E)^2.
\]

\textbf{Paso 3. Expansión y simplificación algebraica.}

Desarrollamos:
\[
r^2 = a^2\cos^2 E - 2a^2e\cos E + a^2e^2 + b^2\sin^2 E.
\]
Sustituimos $b^2 = a^2(1-e^2)$:
\[
r^2 = a^2\cos^2 E - 2a^2e\cos E + a^2e^2 + a^2(1-e^2)\sin^2 E.
\]
Factorizando $a^2$,
\[
r^2 = a^2\left[\cos^2 E + \sin^2 E - e^2\sin^2 E - 2e\cos E + e^2\right].
\]
Usamos la identidad $\cos^2 E + \sin^2 E = 1$:
\[
r^2 = a^2\left[1 - e^2\sin^2 E - 2e\cos E + e^2\right].
\]
Ahora agrupamos los términos con $e^2$ y empleamos $1-\sin^2 E = \cos^2 E$:
\[
r^2 = a^2\left[1 - 2e\cos E + e^2(1-\sin^2 E)\right]
= a^2\left[1 - 2e\cos E + e^2\cos^2 E\right].
\]
Reconocemos un cuadrado perfecto:
\[
r^2 = a^2(1-e\cos E)^2.
\]

\textbf{Paso 4. Extracción de la raíz.}

Como $0 \le e < 1$, el factor $1-e\cos E$ es positivo. Entonces,
\[
\boxed{r = a(1-e\cos E)}.
\]
Esta expresión muestra que la distancia radial toma su mínimo en el periapsis, $r_{\min}=a(1-e)$, y su máximo en el apoapsis, $r_{\max}=a(1+e)$.

\textbf{Inciso b) Demostración de $\tan\frac{f}{2} = \sqrt{\frac{1+e}{1-e}}\tan\frac{E}{2}$.}

\textbf{Paso 1. Igualación de expresiones radiales.}

La ecuación polar de la elipse en función de la anomalía verdadera $f$ es
\[
r = \frac{a(1-e^2)}{1+e\cos f}.
\]
Igualamos esta expresión con la obtenida en el inciso anterior:
\[
a(1-e\cos E) = \frac{a(1-e^2)}{1+e\cos f}.
\]
Cancelando $a$ y despejando,
\[
1+e\cos f = \frac{1-e^2}{1-e\cos E}.
\]

\textbf{Paso 2. Despeje de $\cos f$.}

Restamos $1$ y simplificamos:
\[
e\cos f = \frac{1-e^2}{1-e\cos E} - 1
= \frac{1-e^2-(1-e\cos E)}{1-e\cos E}
= \frac{e(\cos E-e)}{1-e\cos E}.
\]
Dividiendo entre $e$,
\[
\cos f = \frac{\cos E-e}{1-e\cos E}.
\]

\textbf{Paso 3. Aplicación de la identidad del ángulo mitad.}

Usamos
\[
\tan^2\frac{f}{2} = \frac{1-\cos f}{1+\cos f}.
\]
Sustituyendo la expresión anterior para $\cos f$,
\[
\tan^2\frac{f}{2} = \frac{1-\dfrac{\cos E-e}{1-e\cos E}}{1+\dfrac{\cos E-e}{1-e\cos E}}.
\]

\textbf{Paso 4. Simplificación de numerador y denominador.}

Para el numerador,
\[
1-\cos f = \frac{(1-e\cos E) - (\cos E-e)}{1-e\cos E}
= \frac{(1+e)(1-\cos E)}{1-e\cos E}.
\]
Para el denominador,
\[
1+\cos f = \frac{(1-e\cos E) + (\cos E-e)}{1-e\cos E}
= \frac{(1-e)(1+\cos E)}{1-e\cos E}.
\]
Por tanto,
\[
\tan^2\frac{f}{2}
= \frac{\dfrac{(1+e)(1-\cos E)}{1-e\cos E}}{\dfrac{(1-e)(1+\cos E)}{1-e\cos E}}
= \frac{1+e}{1-e}\cdot\frac{1-\cos E}{1+\cos E}.
\]

\textbf{Paso 5. Reconocimiento de la identidad para $E/2$.}

Observamos que
\[
\frac{1-\cos E}{1+\cos E} = \tan^2\frac{E}{2}.
\]
Así,
\[
\tan^2\frac{f}{2} = \frac{1+e}{1-e}\tan^2\frac{E}{2}.
\]
Tomando la raíz positiva,
\[
\boxed{\tan\frac{f}{2} = \sqrt{\frac{1+e}{1-e}}\tan\frac{E}{2}}.
\]

\textbf{Interpretación analítica.}

Esta relación permite convertir el ángulo auxiliar $E$, medido desde el centro geométrico, en la anomalía verdadera $f$, medida desde el foco gravitacional. El factor
\[
\sqrt{\frac{1+e}{1-e}} > 1
\]
para toda elipse no circular, lo que refleja la distinta rapidez angular del cuerpo a lo largo de la órbita, en concordancia con la segunda ley de Kepler.
\end{solution}

\problemsection{Problema 10: Elementos orbitales}{
Para un satélite en órbita alrededor de la Tierra se tienen los siguientes elementos orbitales:
\[
a=8016{,}0\ \text{km},
\quad e=0{,}06,
\quad I=50^\circ,
\quad \Omega=0{,}0^\circ,
\quad \omega=30^\circ,
\quad f=20^\circ.
\]
Encuentre el vector de posición en el plano fundamental del ecuador terrestre.
}

\begin{solution}
\textbf{Paso 1. Cálculo de la distancia radial $r$.}

La ecuación polar de la trayectoria relativa en coordenadas perifocales es
\[
r = \frac{p}{1+e\cos f},
\]
donde el semilatus rectum se calcula mediante
\[
p = a(1-e^2).
\]
Sustituyendo los valores dados,
\[
p = 8016{,}0(1-0{,}06^2) = 8016{,}0(0{,}9964) = 7987{,}1424\ \text{km}.
\]
Entonces,
\[
r = \frac{7987{,}1424}{1+0{,}06\cos 20^\circ}
= \frac{7987{,}1424}{1+0{,}06(0{,}9396926)}
\approx 7560{,}84\ \text{km}.
\]

\textbf{Paso 2. Vector posición en el sistema perifocal.}

En el sistema perifocal $\mathcal{F}_{PQW}$, con el eje $\hat{P}$ dirigido al periapsis y el eje $\hat{Q}$ contenido en el plano orbital, el vector posición es
\[
\vec{r}_{PQW} =
\begin{pmatrix}
r\cos f \\
r\sin f \\
0
\end{pmatrix}.
\]

\textbf{Paso 3. Transformación al sistema ecuatorial terrestre.}

Para pasar del sistema perifocal al sistema ecuatorial inercial se usa la rotación de Euler
\[
\mathbf{R}_{ECI \leftarrow PQW} = \mathbf{R}_z(\Omega)\mathbf{R}_x(I)\mathbf{R}_z(\omega).
\]
De esta transformación se obtienen las componentes cartesianas
\[
\begin{cases}
x = r\left[\cos\Omega\cos(\omega+f) - \cos I\sin\Omega\sin(\omega+f)\right], \\
y = r\left[\sin\Omega\cos(\omega+f) + \cos I\cos\Omega\sin(\omega+f)\right], \\
z = r\left[\sin I\sin(\omega+f)\right].
\end{cases}
\]
Como $\Omega=0^\circ$, se tiene $\cos\Omega=1$ y $\sin\Omega=0$, de modo que las expresiones se simplifican a
\[
\begin{cases}
x = r\cos(\omega+f), \\
y = r\cos I\sin(\omega+f), \\
z = r\sin I\sin(\omega+f).
\end{cases}
\]

\textbf{Paso 4. Sustitución numérica.}

Primero calculamos el ángulo compuesto:
\[
\omega+f = 30^\circ + 20^\circ = 50^\circ.
\]
Luego,
\[
\cos 50^\circ \approx 0{,}6427876,
\qquad
\sin 50^\circ \approx 0{,}7660444.
\]
Como además $I=50^\circ$, se tiene también
\[
\cos I = \cos 50^\circ \approx 0{,}6427876,
\qquad
\sin I = \sin 50^\circ \approx 0{,}7660444.
\]
Sustituyendo en las componentes del vector posición,
\[
x = 7560{,}84(0{,}6427876) \approx 4860{,}0\ \text{km},
\]
\[
y = 7560{,}84(0{,}6427876)(0{,}7660444) \approx 3722{,}8\ \text{km},
\]
\[
z = 7560{,}84(0{,}7660444)(0{,}7660444) \approx 4436{,}9\ \text{km}.
\]

\textbf{Resultado final.}

El vector de posición del satélite en el sistema de referencia ecuatorial terrestre es
\[
\boxed{\vec{r} \approx
\begin{pmatrix}
4860{,}0 \\
3722{,}8 \\
4436{,}9
\end{pmatrix}\ \text{km}}.
\]
Este vector tiene magnitud radial aproximada $r\approx 7560{,}8\ \text{km}$ y es consistente con la inclinación orbital de $50^\circ$ respecto del plano ecuatorial.
\end{solution}


\end{document}
